\documentclass[12pt]{article}
\usepackage[utf8]{inputenc}
\usepackage[T1]{fontenc}
\usepackage[margin=1in]{geometry}
\usepackage{setspace}
\usepackage{xcolor}
\usepackage[colorlinks=true, urlcolor=blue!60!black, linkcolor=blue!60!black]{hyperref}
\setlength{\parindent}{0em}
\setlength{\parskip}{0.45em}
\renewcommand{\baselinestretch}{1.0}
\begin{document}
\begin{center}
{\Large \textit{Wisconsin v. Mitchell}, 508 U.S. 476 (1993)\par}
\vspace{0.3em}
{\normalsize Case Brief --- One-Page Edition\par}
\end{center}

\textbf{Memory Jogger.} A hate-crime penalty enhancement for choosing a victim by race punishes conduct, not thought, and survives the First Amendment (\href{https://www.courtlistener.com/opinion/112879/wisconsin-v-mitchell/}{CourtListener}).

\textbf{Facts.} In Kenosha in 1989, after discussing a scene from \textit{Mississippi Burning}, Todd Mitchell urged his group to ``move on some white people'' and sent them after a passing white boy, who was beaten, robbed of his shoes, and left in a coma for four days. A Wisconsin statute raised the maximum penalty by up to five years when a defendant selects a victim because of race or another listed trait.

\textbf{Procedural History.} Convicted of aggravated battery with a racial-selection finding, Mitchell faced seven years instead of two and received four. The Wisconsin Supreme Court held the enhancement punished ``offensive thought.''

\textbf{Votes.} Reversed 9--0 (Rehnquist); no separate opinions.

\textbf{Holding.} A State may enhance the penalty for unprotected conduct when the victim is chosen because of a protected trait, and prior speech may prove motive or intent.

\textbf{Reasoning.} Physical assault is not expressive conduct, and sentencing has long weighed motive; \textit{Dawson v. Delaware} bars only punishing ``abstract beliefs'' unconnected to the crime. Like Title VII, the statute targets discriminatory conduct; unlike the ordinance in \textit{R.A.V.}, it does not target expression. Bias crimes inflict greater harm, and the feared chilling effect was speculative: few would suppress bigoted views for fear of a future felony trial. Evidence of speech remains subject to ordinary rules of relevance.

\textbf{Examples.} \emph{Same-Side:} enhanced punishment for burning a church because of religion. \emph{Opposite-Side:} jailing people for epithets alone. \emph{Fence-Sitter:} a protester who screams slurs and shoves a sign.

\textbf{Critique.} Civil libertarians such as Susan Gellman say the extra years punish only motive and let legislatures evade \textit{R.A.V.}; civil-rights scholars such as Frederick Lawrence answer that bias violence terrorizes whole communities.

\textbf{Key Quotation.} ``The First Amendment, moreover, does not prohibit the evidentiary use of speech to establish the elements of a crime or to prove motive or intent.'' --- Chief Justice Rehnquist, 508 U.S. at 489.

\end{document}
